\documentclass[12pt]{article}
\usepackage[utf8]{inputenc}
\usepackage[spanish]{babel}
\usepackage{float}
\usepackage{graphicx}
\usepackage{fancyhdr}
\usepackage{graphics}
\usepackage[
    type={CC},
    modifier={by-nc-sa},
    version={4.0},
]{doclicense}
\usepackage{mdframed}
\usepackage[margin = 1.2in]{geometry}
\usepackage{listings} 
\usepackage{color}
\usepackage[listings]{tcolorbox}

\definecolor{codegreen}{rgb}{0,0.6,0} \definecolor{codegray}{rgb}{0.5,0.5,0.5} \definecolor{codepurple}{rgb}{0.58,0,0.82} \definecolor{backcolour}{rgb}{0.95,0.95,0.92}

\lstdefinestyle{mystyle}{ backgroundcolor=\color{backcolour}, commentstyle=\color{ codegreen}, 
keywordstyle=\color{magenta}, numberstyle=\tiny\color{codegray}, stringstyle=\color{codepurple}, 
basicstyle=\footnotesize,
breakatwhitespace=false,
breaklines=true,
captionpos=b,
keepspaces=true,
numbers=left,
numbersep=5pt,
showspaces=false,
showstringspaces=false,
showtabs=false,
tabsize=2
}

\lstset{style=mystyle}


\title{El manual básico de CoreWar \\
\large Aprende Redcode DAT cero.}
\author{Fernando Méndez Torrubiano}
\date{2 de febrero de 2019}

\begin{document}
\pagestyle{fancy} %Estilo de página a partir de aquí
\rhead{El manual básico de CoreWar} %Encabezado a la derecha.
\lhead{Fernando Méndez}

\maketitle

\begin{figure}[H]
    \centering
    \includegraphics[width= 0.8\textwidth]{Media/COREWARUCM.jpg}
\end{figure}

\newpage
\tableofcontents

\newpage
\section{Introducción}
Bienvenidos a la \testxtbf{Guía del auto-Hacker Galáctico}, un breve documento en el que se explicarán las nociones básicas para comenzar tu camino como \textit{hacker} en \textbf{FDIst}.
\newline
\noindent
\subsection{Sistemas Operativos}
\noindent
He de comenzar diciendo que todo lo que se explique aquí, irá enfocado a la plataforma \testxtbf{Linux}, sin embargo, si usas \textbf{Windows}, no te preocupes, pues te enseñaremos a crear y usar una máquina virtual para que puedas trabajar sin necesidad de tener ninguna distribución de \textbf{GNU/Linux}.
\newline\\
Además, recalcar que cada uno es libre de usar la plataforma en la que más cómodo se sienta, no por usar \textbf{Windows o MacOS}, serás un informático de segunda, las herramientas están ahí y gracias a internet, toda la documentación está a vuestro alcance. Por tanto, si preferís usar otro Sistema Operativo (SO), que no esté basado en GNU/Linux, sabed que tenéis herramientas equivalentes a las que se van a detallar en esta guía. Eso sí, \textbf{tened cuidado con qué instaláis}, pues existen programas de dudosa seguridad y que incluyen spyware.
\subsection{Advertencia y VPN}
\noindent
Por último advertíos de que \textbf{FDIst} tiene un acuerdo con el \textbf{CERN}, el \textbf{CSN} y el \textbf{Centro de Procesamiento de Datos (CPD)} de la Universidad Complutense de Madrid (UCM). Cualquier ataque a otra entidad o empresa, no estará respaldado por FDIst y \textbf{todo ataque que se realice a los centros anteriormente mencionados, deberá ser a través de la VPN} que se os proporcionará en este mismo manual y a la que previamente se os habrá dado acceso a través de vuestro correo UCM.
\newpage
\section{Instrucciones básicas.}
%\noindent

\begin{tcolorbox}
    \begin{itemize}
        \item \textbf{LAS DIRECCIONES DE MEMORIA (\texttt{DIRECCIONAMIENTO DIRECTO}) SE ESCRIBEN CON '\$' [O DIRECTAMENTE CON EL ENTERO]} (se da una dirección de memoria que contiene un valor).
        \item \textbf{LOS INMEDIATOS CON UNA '\#' DELANTE DEL \newline NÚMERO} (representa el valor del entero dado).
        \item \textbf{LOS COMENTARIOS SE ESCRIBEN CON UN ';'.}
    \end{itemize}
\end{tcolorbox}

\subsection{DAT}
\noindent
\textbf{Mata el proceso}. 
\begin{itemize}
    \item Ejemplo:
    \begin{itemize}
        \item \textbf{\textsc{DAT \#0, \#0}}
    \end{itemize}
\end{itemize}

\newline (Casi imprescindible si queremos ganar, aunque se puede volver en nuestra contra).

\subsection{MOV}
\noindent
\textbf{Copia el dato de una dirección a otra}.
\begin{itemize}
    \item Ejemplos:
    \begin{itemize}
        \item \textbf{\textsc{MOV dir1, dir2}} ;Copia el contenido de la dirección 'dir1' en la 'dir2'.
        \item \textbf{\textsc{MOV inm, dir}} ;Copia el valor del inmediato 'inm' en la dirección 'dir'.
    \end{itemize}
\end{itemize}


\subsection{ADD/SUB}
\noindent
\textbf{Suma / Resta}.
\begin{itemize}
    \item Ejemplos:
    \begin{itemize}
        \item \textbf{\textsc{ADD dir1, dir2}} ;Suma el contenido de la dirección 'dir1' al de la 'dir2'.
        \item \textbf{\textsc{ADD inm, dir}} ;Suma el valor 'inm' al contenido de la dirección 'dir'.
        \item \textbf{\textsc{SUB dir1, dir2}} ;Resta el contenido de la dirección 'dir1' al de la 'dir2'.
        \item \textbf{\textsc{SUB inm, dir}} ;Resta el valor 'inm' al contenido de la dirección 'dir'.
    \end{itemize}
\end{itemize}

\subsection{JMP}
\noindent
\textbf{Salto incondicional}.
\begin{itemize}
    \item Ejemplos:
    \begin{itemize}
        \item \textbf{\textsc{JMP dir}} ;Salta a la dirección 'dir'.
        \item \textbf{\textsc{JMP etiqueta}} ;Salta a la dirección de la etiqueta.
    \end{itemize}
\end{itemize}
\newpage
\section{Instrucciones avanzadas.}
%\noindent

\begin{tcolorbox}
    \begin{itemize}
        \item \textbf{LOS DIRECCIONAMIENTOS DIRECTOS SE ESCRIBEN CON '\$' \newline [O DIRECTAMENTE CON EL ENTERO]} \newline (se da una dirección de memoria que contiene un valor).
        \item \textbf{LOS DIRECCIONAMIENTOS INDIRECTOS SE ESCRIBEN CON '@'} \newline (se da una dirección de memoria que contiene otra dirección de memoria, la cual contiene, ahora sí, el valor).
        \item \textbf{LOS DIRECCIONAMIENTOS INDIRECTOS CON \newline PREDECREMENTO SE ESCRIBEN CON '$<$'}
        \item \textbf{LOS DIRECCIONAMIENTOS INDIRECTOS CON \newline POSTINCREMENTO SE ESCRIBEN CON '$>$'}
        \item \textbf{EXISTEN CONSTANTES COMO EL TAMAÑO DEL NÚCLEO (\texttt{CORESIZE}) O EL NÚMERO DE PROGRAMAS EN EJECUCIÓN (\texttt{WARRIORS}).}
    \end{itemize}
\end{tcolorbox}

\subsection{MUL/DIV}
\textbf{Multiplicación / División}.
\begin{itemize}
    \item Ejemplos:
    \begin{itemize}
        \item \textbf{MUL dir1, dir2} ;Multiplica el contenido de la dirección 'dir1' por el contenido de la dirección 'dir2'.
        \item \textbf{DIV imd1, imd2} ;Divide el número 'imd1' entre el número 'imd2'.
    \end{itemize}
\end{itemize}

\subsection{JMZ}
\textbf{Salto condicional. Comprueba un número y salta si es cero}.
\begin{itemize}
    \item Ejemplos:
    \begin{itemize}
        \item \textbf{JMZ dir1, dir2} ;Salta a la dirección 'dir1' si 'dir2' es cero.
    \end{itemize}
\end{itemize}

\subsection{JMN}
\textbf{Salto condicional. Comprueba un número y salta si NO es cero}.
\begin{itemize}
    \item Ejemplos:
    \begin{itemize}
        \item \textbf{JMN dir1, dir2} ;Salta a la dirección 'dir1' si 'dir2' NO es cero.
    \end{itemize}
\end{itemize}

\subsection{DJN}
\textbf{Resta uno a un número y salta si el resultado de la resta NO es 0}.
\begin{itemize}
    \item Ejemplos:
    \begin{itemize}
        \item \textbf{DJN dir1, dir2} ;Decrementa uno al contenido de 'dir1', si el resultado NO es cero, salta a la dirección 'dir2'.
    \end{itemize}
\end{itemize}

\subsection{SPL}
\textbf{Crea un nuevo proceso en otra dirección}.
\begin{itemize}
    \item Ejemplos:
    \begin{itemize}
        \item \textbf{SPL dir} ;Subdivisión del programa. Comienza un nuevo proceso en la dirección 'dir'.
        \item \textbf{SPL etiqueta} ;Subdivisión del programa, añadiéndose al proceso o procesos en ejecución el situado en la dirección de la 'etiqueta'.
    \end{itemize}
\end{itemize}

\subsection{CMP/SEQ}
\textbf{Salto condicional. Compara dos números, si son iguales, se salta la siguiente instrucción}.
\begin{itemize}
    \item Ejemplos:
    \begin{itemize}
        \item \textbf{SEQ dir1, dir2} ;Si el contenido de la dirección 'dir1', es igual al de 'dir2', no ejecuta la siguiente instrucción.
    \end{itemize}
\end{itemize}

\subsection{SNE}
\textbf{Salto condicional. Compara dos números, si NO son iguales, se salta la siguiente instrucción}.
\begin{itemize}
    \item Ejemplos:
    \begin{itemize}
        \item \textbf{SNE dir1, dir2} ;Si el contenido de la dirección 'dir1', NO es igual al de 'dir2', no ejecuta la siguiente instrucción.
    \end{itemize}
\end{itemize}

\subsection{NOP}
\textbf{No hace nada}.
\begin{itemize}
    \item Ejemplos:
    \begin{itemize}
        \item \textbf{NOP dir1, dir2} ;Pues eso, no hace nada. Pero los operandos se siguen evaluando, es decir, consume un ciclo de ejecución.
    \end{itemize}
\end{itemize}
\newpage
\section{Mi primer virus.}
\noindent
Lo primero es saber qué instrucciones están dentro de nuestro alcance y cuales sobrepasan nuestros conocimientos, es decir, hay que ser realistas, si es la primera vez que programas en un lenguaje ensamblador no es necesario que hagas un virus con decenas de saltos y líneas de código.
\newline Y es tan sencillo como ver los ejemplos, entenderlos, coger uno que nos guste y mejorarlo. 
\begin{flushright} Fernando \end{flushright}

\subsection{Ejemplo de ejecución instrucción a instrucción}
\begin{tcolorbox}
    ADD \#4, 3 \phantom{................};Aquí comienza la ejecución. \\
    MOV 2, @2\\
    JMP -2\\
    DAT \#0, \#0
\end{tcolorbox}

\begin{tcolorbox}
    ADD \#\textbf{4}, 3\\
    MOV 2, @2 \phantom{................};Siguiente instrucción. \\
    JMP -2\\
    DAT \#0, \#\textbf{4}
\end{tcolorbox}

\begin{tcolorbox}
    ADD \#4, 3\\
    MOV \textbf{2}, \textbf{@2} ;\phantom{...}---------,\\
    JMP -2\phantom{........} ;\phantom{.............}\| \textbf{+2}\\
    DAT \#0, \#\textbf{4} ;\phantom{..}<--' --,\\ %El operando B del MOV, apunta aquí.\\
    ...\phantom{.......................................} \|\\
    ...\phantom{.......................................} \| \textbf{+4}\\
    ...\phantom{.......................................} \|\\
    DAT \#0, \#4 ;\phantom{...}<----' %El operando B del DAT apunta aquí.
\end{tcolorbox}

\begin{tcolorbox}
    ADD \#4, 3\\
    MOV 2, @2\\
    JMP -2\phantom{................};Siguiente instrucción. \\
    DAT \#0, \#8\\
    ...\\
    ...\\
    ...\\
    DAT \#0, \#4\\
    ...\\
    ...\\
    ...\\
    DAT \#0, \#8
\end{tcolorbox}

\begin{tcolorbox}
    ADD \#4, 3\phantom{................};Siguiente instrucción. \\
    MOV 2, @2\\
    JMP -2\\
    DAT \#0, \#8\\
    ...\\
    ...\\
    ...\\
    DAT \#0, \#4\\
    ...\\
    ...\\
    ...\\
    DAT \#0, \#8
\end{tcolorbox}
\newpage
\section{Estrategias.}
\noindent
\textsc{Siempre se halla la eterna duda entre \textbf{FUERZA} y \textbf{VELOCIDAD}.
\newline Rellenar rápidamente el núcleo o atacar a posiciones estratégicas.} \\
\\ En CoreWar existen \textbf{3 estrategias} fundamentales que, por analogía con el famoso juego, se les denominan \textbf{piedra, papel y tijera}.

\begin{tcolorbox}[colframe=black!75!blue]
    \begin{enumerate}
        \item \textbf{PAPEL: } \underline{hace múltiples copias de sí mismo lo más rápidamente posible}, así sacrifica velocidad de ataque por resistencia.
        \newline Esta estrategia \underline{vence a piedra} pero \underline{pierde ante tijeras} gracias a su gran capacidad de supervivencia aunque tienen una cierta tendencia al empate. 
        \item  \textbf{PIEDRA: } \underline{bombardea direcciones de memoria a ciegas} intentando matar rápidamente al mayor número de enemigos. Su reducido tamaño y sencillez los hace relativamente robustos y difíciles de localizar.
        \newline Esta estrategia \underline{vence a tijeras} pero \underline{pierde ante papel}. 
         \item \textbf{TIJERA: } \underline{es la más avanzada.} Comprueba posiciones de memoria a intervalos hasta localizar al guerrero rival. Una vez localizado generalmente sobrescriben su código con instrucciones que les obligan a generar nuevos procesos indefinidamente hasta quedar prácticamente bloqueados. Después proceden a eliminar todos los rivales.
         \newline Esta estrategia generalmente \underline{vence al papel} y \underline{pierde contra piedra}, puesto que pierde tiempo atacando las posiciones de memoria alteradas por este último. 
    \end{enumerate}
\end{tcolorbox}
\newpage
\section{Ejemplos.}
\noindent

\begin{enumerate}
    \item \underline{TRASGO}: \\
    \newline \textbf{MOV 0, 1} ;Copia el contenido de la dirección '0' en la dirección '1'. \newline \phantom{....................};No puede ganar, sólo empatar.
    \newline

    \item \underline{ENANO BOMBARDERO} (DRAW): \\
    \newline \textbf{ADD \#4, 3} \phantom{...};Suma al contenido de la dirección '4', el entero '3'.
    \newline \textbf{MOV 2, @2} \phantom{...};Copia el contenido de la dirección '2' en la dirección contenida \newline \phantom{..........................}en la dirección '2'.
    \newline \textbf{JMP -2} \phantom{...........};Salta a la dirección de memoria '-2'.
    \newline \textbf{DAT \#0, \#0} \phantom{.};Termina el proceso.
    \newline
    
    \item \underline{BLANCA-!NIEVES:}
    \begin{tcolorbox}
        \newline \textit{(El siguiente código es parte del virus 'BLANCA-NIEVES', observando atentamente, veremos que son una serie de 'ENANOS' encadenados de una forma sofisticada).}
    \end{tcolorbox}
    \newline \textbf{Enano: \phantom{.....}SPL Mudito}
    \newline \textbf{Sabiondo: ADD \#1328, 3}       
    \newline \textbf{\phantom{.................}MOV 2, @2}
    \newline \textbf{\phantom{.................}JMP -2}
    \newline \textbf{\phantom{.................}DAT \#0, \#0}
    \newline \textbf{Mudito: \phantom{...}ADD \#516,3}    
    \newline \textbf{\phantom{.................}MOV 2, @2}
    \newline \textbf{\phantom{.................}JMP -2} 
    \newline \textbf{\phantom{.................}DAT \#0, \#0}
    \newline
    
    \item \underline{ZERG-RUSH:} \\
    \newline \textbf{salt\phantom{..}EQU \#100} \phantom{.......};Etiqueta 'salt'. Establece una constante, cuyo valor será \phantom{.....................................}100.
    \newline \textbf{ \phantom{.......}MOV salt, 10} \phantom{...};Copia el valor de la dirección de la etiqueta 'salt' en la\newline \phantom{......................................}de direcciónmemoria '10'.
    \newline \textbf{loop MOV imp, @9} \phantom{.};Etiqueta 'loop'.
	\newline \textbf{\phantom{........}SPL @8} \phantom{.............};Crea un nuevo proceso en la dirección que contiene \newline \phantom{......................................}la dirección de memoria '8'.
	\newline \textbf{\phantom{........}ADD \#100, 7} \phantom{..};Suma '100' al contenido de la dirección '7'
	\newline \textbf{\phantom{........}JMP loop} \phantom{.........};Salta a la etiqueta 'loop'
    \newline \textbf{imp \phantom{.}MOV 0,1} \phantom{..........};Etiqueta 'imp'. Copia el contenido de la dirección '0' en \phantom{.....................................} la dirección '1. \\
    \newline \phantom{....................................................}(Cortesía del Profesor José Luis Vázquez-Poletti).
    
\end{enumerate}
\newpage
\vspace*{\fill} % Para que lo escriba al final de la página
\doclicenseThis %done

\end{document}
